\begin{example}
    Consideremos una sucesión de racionales $r=(r_n)_{n\geq1}$. Definimos la sucesión 
    $x=(x_n)_{n\geq1}$ por
    \[
        x_n = 
        \begin{cases} 
            r_{n/2} & \text{si $n$ es par} \\ 
            0 & \text{si $n$ es impar}
        \end{cases} 
    \]

    Entonces $(x_n)$ no es equidistribuida módulo $1$.
\end{example}

\begin{resolve}
    Consideramos $[0,b) \subset \mathbb{T}$ con $b < 1/2$. Si la sucesión fuese
    equidistribuida módulo $1$, debería cumplirse que
    \[
        \lim_{N\to\infty} \frac{A_N([0,b),x)}{N}  = b
    \]

    Sin embargo, para $N=2M$,
    \[
        A_{2M}([0,b),x) = M + R \qquad \text{donde } R = A_M([0,b),r) 
    \]
    Los índices impares de la sucesión $(x_n)$ son iguales a $0$ y contribuyen
    con $M$ términos.
    
    Por consiguiente,
    \[
        \frac{A_{2M}([0,b),x)}{2M} = 
        \frac{M+R}{2M} \geq  
        \frac{M}{2M} = 
        \frac{1}{2}
    \] 

    Debido a que $b < 1/2$, el límite anterior no puede ser igual a $b$ 
    y por tanto $(x_n)$ no es equidistribuida módulo $1$.
\end{resolve}